\PassOptionsToPackage{no-math}{fontspec}
\documentclass[UTF8,zihao=-4,openany,twoside,fontset=fandol]{ctexbook}
\usepackage[paperwidth=190mm,paperheight=270mm,inner=23mm,outer=19mm,top=24mm,bottom=25mm,headheight=15pt,headsep=8mm,marginparwidth=13mm,marginparsep=3mm]{geometry}
\usepackage{fontspec}
\setmainfont{texgyrepagella}[Extension=.otf,UprightFont=*-regular,BoldFont=*-bold,ItalicFont=*-italic,BoldItalicFont=*-bolditalic]
\setsansfont{texgyreheros}[Extension=.otf,UprightFont=*-regular,BoldFont=*-bold,ItalicFont=*-italic,BoldItalicFont=*-bolditalic,Scale=MatchLowercase]
\usepackage{amsmath,amssymb,amsthm,mathtools}
\usepackage{newpxmath}
\usepackage{graphicx,xcolor,tikz,tikz-cd,booktabs,longtable,array,enumitem,microtype,fancyhdr,needspace,etoolbox}
\usetikzlibrary{arrows.meta,calc,decorations.pathreplacing,positioning,matrix}
\definecolor{Ink}{HTML}{173947}
\definecolor{Accent}{HTML}{A5844A}
\definecolor{Muted}{HTML}{69777C}
\definecolor{Paper}{HTML}{F7F5EF}
\usepackage[unicode,colorlinks,linkcolor=Ink,citecolor=Ink,urlcolor=Ink,bookmarksnumbered=true,pdfdisplaydoctitle=true]{hyperref}
\hypersetup{pdftitle={算术代数几何：C.I.M.E. 1991年讲义（中文译本）},pdfauthor={Jean-Louis Colliot-Thélène; Kazuya Kato; Paul Vojta},pdfsubject={Lecture Notes in Mathematics 1553 中文翻译，附29项译者补注}}
\urlstyle{same}
\ctexset{chapter={format=\Huge\bfseries\color{Ink},name={第,篇},number=\chinese{chapter},beforeskip=10pt,afterskip=24pt},section={format=\Large\bfseries\color{Ink},beforeskip=18pt,afterskip=9pt},subsection={format=\large\bfseries\color{Ink},beforeskip=14pt,afterskip=7pt}}
\setlength{\parindent}{2em}
\linespread{1.16}
\setlength{\emergencystretch}{2em}
\allowdisplaybreaks[2]
\setlist[itemize]{leftmargin=2em,itemsep=3pt,topsep=5pt}
\setlist[enumerate]{leftmargin=2.4em,itemsep=3pt,topsep=5pt}
\setlist[description]{style=nextline,font=\normalfont\bfseries,leftmargin=1.7em,itemsep=5pt}
\pagestyle{fancy}
\fancyhf{}
\fancyhead[LE]{\small\color{Muted}算术代数几何}
\fancyhead[RO]{\small\color{Muted}\begingroup\termanchorsfalse\nouppercase{\leftmark}\endgroup}
\fancyfoot[LE,RO]{\color{Ink}\thepage}
\renewcommand{\headrulewidth}{0.3pt}
\renewcommand{\chaptermark}[1]{\markboth{\ifcase\value{chapter}#1\or 挠闭链与代数 $K$ 理论\or Hasse--Weil $L$ 函数与岩泽理论\or 丢番图逼近\else #1\fi}{}}
\fancypagestyle{plain}{\fancyhf{}\fancyfoot[LE,RO]{\color{Ink}\thepage}\renewcommand{\headrulewidth}{0pt}}
\newcommand{\chapterauthor}[1]{\begin{flushright}\large\color{Muted}#1\end{flushright}\vspace{1.5em}}
\newcommand{\origpage}[2]{\par\phantomsection\label{orig:#1}\marginpar{\raggedright\scriptsize\color{Muted}原书\,#2}}
\newif\iftermanchors
\termanchorstrue
\DeclareRobustCommand{\termuse}[4]{\iftermanchors\hypertarget{termocc-#2}{\label{termocc-#2}\hyperlink{term-#1}{#3（#4）}}\else\hyperlink{term-#1}{#3（#4）}\fi}
\DeclareRobustCommand{\term}[3]{#1（#2）}
\newenvironment{statement}[1]{\par\addvspace{8pt}\needspace{3\baselineskip}\noindent\textbf{#1}\quad\ignorespaces}{\par\addvspace{7pt}}
\renewcommand{\proofname}{证明}
\newcommand{\Spec}{\operatorname{Spec}}
\newcommand{\Gal}{\operatorname{Gal}}
\newcommand{\Pic}{\operatorname{Pic}}
\newcommand{\Br}{\operatorname{Br}}
\newcommand{\CH}{\operatorname{CH}}
\newcommand{\Hom}{\operatorname{Hom}}
\newcommand{\Ext}{\operatorname{Ext}}
\newcommand{\Tor}{\operatorname{Tor}}
\newcommand{\coker}{\operatorname{Coker}}
\newcommand{\Ker}{\operatorname{Ker}}
\newcommand{\im}{\operatorname{Im}}
\newcommand{\rank}{\operatorname{rank}}
\newcommand{\NS}{\operatorname{NS}}
\newcommand{\Alb}{\operatorname{Alb}}
\newcommand{\et}{\mathrm{\acute et}}
\newcommand{\dR}{\mathrm{dR}}
\newcommand{\cris}{\mathrm{cris}}
\newcommand{\tors}{\mathrm{tors}}
\newcommand{\id}{\operatorname{id}}
\newcommand{\Tr}{\operatorname{Tr}}
\newcommand{\Nm}{\operatorname{Nm}}
\pdfstringdefDisableCommands{\def\term#1#2#3{#1}\def\termuse#1#2#3#4{#3}\def\mathbb#1{#1}\def\mathcal#1{#1}\def\dR{dR}}

% 补注编号与2026-11-05收到的勘误说明一致。
\newcommand{\DeclareTranslatorNote}[3]{%
  \csgdef{translator-note-title-#1}{#2}%
  \csgdef{translator-note-text-#1}{#3}}
\DeclareRobustCommand{\translatornote}[1]{%
  \footnote{\hypertarget{translator-note-#1}{}\label{translator-note:#1}%
    \textbf{译者补注\,\hyperlink{translator-index-#1}{#1}。}%
    \csuse{translator-note-text-#1}}}
\DeclareRobustCommand{\translatornoteref}[1]{%
  \textsuperscript{\hyperlink{translator-note-#1}{〔补注#1〕}}}
% 依据用户于2026-11-05提供的勘误说明编写；编号与该说明一一对应。
% 本文件由主文件在导言区载入。正文保留原式，在相应位置调用译者补注。
\DeclareTranslatorNote{1}{Picard 下降与可分闭包}{这里的下降正合列应采用可分闭包 $k^s$，即
\[
0\longrightarrow\Pic(X)\longrightarrow\Pic(X_{k^s})^{\Gal(k^s/k)}\longrightarrow\Br(k).
\]
若坚持使用代数闭包，则须增加 $k$ 完美的假设；$X(k)\ne\varnothing$ 时的同构断言也应作同样处理。例如，取 $k=\mathbb F_2(t)^{\mathrm{sep}}$ 和 $E:y^2+y=x^3$。此时 $k$ 可分闭，$G=1$ 且 $\Br(k)=0$，但 $\sqrt t\notin k$；在 $\bar k$ 上可取 $P=(\sqrt t,y)$，于是 $\mathcal O(P-O)$ 不能下降到 $k$。因此代数闭包上的 Picard 群不能由所写的不动部分下降式给出。命题 1.1 本身的拉回单射结论仍可保留。}

\DeclareTranslatorNote{2}{定理 6.4 的正特征证明缺口}{本定理的陈述包含正特征基域，但下列证明没有覆盖其全部范围：a) 所用的平展上同调单射，先前只在 $\ell\ne\operatorname{char}k$ 的初级挠元上建立；b) 对剩余核的指数有界论证又明确假设特征为零。因此，正特征情形，特别是 $p$ 初级挠元，不能由这里的证明推出。若仅采用本书给出的论证，应把 $k$ 限为 $\mathbb Q$ 上有限生成的域；保留原陈述则需另补正特征证明。这是证明范围的缺口，不能只靠改动一个符号补足。}

\DeclareTranslatorNote{3}{Brauer 群目标应取模常数类的商}{这里的自然对应应以商群为目标。在本段假设下，Hochschild--Serre 谱序列给出
\[
H^1(k,\Pic\bar X)\simeq\Br_1(X)/\im\Br(k),
\]
其中 $\Br_1(X)=\Ker[\Br(X)\to\Br(\bar X)]$。没有另行指定提升所需的结构，不能自然地把这一商群提升到 $\Br(X)$。因此所写复合应在上述商群上解释，相应交换图也应据此取商。常数 Brauer 类与次数为零的闭链群 $A_0(X)$ 的配对为零，故所需配对可自然下降到这个商群。}

\DeclareTranslatorNote{4}{定理 9.1 中的 $H^0/H^2$}{第二个交换图及紧接说明中的两处 $H^0(\bar C,\mathbb Q/\mathbb Z(2))^G$，均应读为 $H^2_{\et}(\bar C,\mathbb Q/\mathbb Z(2))^G$。这里用的是谱序列微分 $d_2^{0,2}:E_2^{0,2}\to E_2^{2,1}$，而曲线的迹同构给出
\[
H^2_{\et}(\bar C,\mathbb Q/\mathbb Z(2))\simeq\mathbb Q/\mathbb Z(1).
\]
真正的 $H^0$ 则是 $\mathbb Q/\mathbb Z(2)$，其 Tate 扭转次数不同，不能由此认作 $\mu(k)$。例如在 $k=\mathbb Q$ 时，错误表达式的不动部分含三阶元，而 $\mu(\mathbb Q)=\{\pm1\}$。}

\DeclareTranslatorNote{5}{$p$ 处 Euler 因子的 Frobenius}{在 $\ell=p$ 的分支，若使用第二章 §1.3 中标准的 Fontaine Frobenius，则应写
\[
\det_\Lambda(1-\varphi_p t;D_{\mathrm{crys}}(V)),
\]
不应把 $\varphi_p$ 取逆。例如 $V=\mathbb Q_p(1)$ 时，$D_{\mathrm{crys}}(V)$ 上的 Frobenius 为 $p^{-1}$，因子应为 $1-p^{-1}t$，现有写法却给出 $1-pt$。这一修正也与本节的 $L(V(r),s)=L(V,s+r)$ 相容。$\ell\ne p$ 分支中的 $\sigma_\ell^{-1}$ 保持原约定。}

\DeclareTranslatorNote{6}{$V$ 的定义缺少 Tate 扭转}{此处应补上 $(r)$：
\[
V=H^m_{\et}(X\otimes_{\mathbb Q}\overline{\mathbb Q},\mathbb Q_p)(r).
\]
紧接着的例子取 $X=\Spec K$、$(m,r)=(0,1)$，并使用 $\mathbb Q_p(1)$ 系数的上同调；原定义不带扭转，只会给出 $\mathbb Q_p$。例如 $K=\mathbb Q$ 时，两种系数的零次上同调分别为 $\mathbb Q_p$ 与 $0$。猜想 2.2.7 的定义已写出所需的 $(r)$。}

\DeclareTranslatorNote{7}{猜想 2.2.7 的例外参数}{条件 (ii) 中第一组例外应为 $(m,r)=(0,1)$，即改为
\[
(m,r)\ne(0,1),(2\dim X,\dim X+1).
\]
这也关系到 2.2.6(4) 中 $H^1_S\simeq H^1_{\mathrm{lim}}$ 的断言。按补注 6 补上 Tate 扭转后，取 $X=\mathbb P^1_{\mathbb Q}$、$m=0,r=1$，原写法仍允许该例，但
\[
H^1_S\simeq\mathbb Z[1/S]^\times\otimes\mathbb Q_p,
\qquad H^1_{\mathrm{lim}}\simeq\mathbb Q^\times\otimes\mathbb Q_p.
\]
前者有限维，后者无限维，不可能同构。因此这里必须排除 $(0,1)$。}

\DeclareTranslatorNote{8}{Lichtenbaum 公式中的 $S$ 因子}{右侧取 $\mathcal O_K[1/S]$ 上的上同调时，左侧应为去掉 $S$ 处 Euler 因子的 $\zeta_K^{(S)}(1-r)_{(p)}$，沿用前文对部分 zeta 函数的定义。若保留完整 $\zeta_K$ 的写法，则须相应限制右侧的 $S$，例如取 $S=\{p\}$。取 $K=\mathbb Q,p=3,r=2$，把 $7$ 加入 $S$ 会引入因子 $1-7=-6$，其 $3$ 幂部分为 $3$，定位正合列使右侧增加同样的因子，而完整 zeta 值不随 $S$ 改变。}

\DeclareTranslatorNote{9}{分圆单位的素数幂判据}{对 $N\ge2$，此处“当且仅当 $N$ 为合数”应为“当且仅当 $N$ 不是素数幂”。例如 $N=4$ 虽为合数，但 $c_4=(1-i)(1+i)=2$，在 $\mathbb Z[i]^+=\mathbb Z$ 中不是单位。正确判据由分圆多项式 $\Phi_N(1)$ 的取值立即得到；后一句关于在 $\mathbb Z[\alpha,1/N]^+$ 中为单位的叙述不受这一端点判据修正的影响。}

\DeclareTranslatorNote{10}{有限域 zeta 函数的零极点阶符号}{这里以零点的阶为正、极点的阶为负，则两处 $(-1)^i$ 都应改为 $(-1)^{i+1}$：
\[
\operatorname{ord}_{t=1}Z_\Lambda(X,D_\Lambda\mathcal F,t)
=\sum_i(-1)^{i+1}\rank_\Lambda P^i
=\sum_i(-1)^{i+1}\rank_\Lambda Q^i.
\]
例如 $X=\Spec\mathbb F_\ell$、$\Lambda=\mathbb Q_p$、$\mathcal F=\mathbb Q_p$ 时，$Z=(1-t)^{-1}$ 的阶为 $-1$，而 $P^0=Q^0=\mathbb Q_p$，其余为零，原右侧却为 $+1$。使用后面的首项归一化时也应同步留意此符号。}

\DeclareTranslatorNote{11}{Frobenius 同余条件的适用范围}{按所收勘误说明，同余条件应只施加于系数环：对 $a\in\mathcal O_K$ 要求 $\phi(a)\equiv a^q\pmod\pi$，再令 $\phi$ 固定形式参数，从而扩张到 $\mathcal O_K[[t]]$。原文同时要求固定 $\mathcal O(\Gamma)$，并对所有函数 $f$ 满足 $\phi(f)\equiv f^q\pmod\pi$；取形式参数 $t\in\mathcal O(\Gamma)$，便同时得到 $\phi(t)=t$ 和 $t\equiv t^q\pmod\pi$，在剩余域上的形式幂级数环中矛盾。后文使用的 $\phi^{-n}$ 应在修正后的系数 Frobenius 约定下理解。}

\DeclareTranslatorNote{12}{微分迹与拉回复合的标量}{此复合应为乘以 $[\pi]$ 的次数，即
\[
\Tr_\pi\circ[\pi]^*=q\,\id_{\Omega^1(G)},
\qquad q=\#(\mathcal O_F/\pi),
\]
而非乘以 $\pi$。这与 2.3.1 后所定义的迹及投影公式相符；在相邻层 $m=n+1$ 取 $x=1$ 即得到标量 $q$。例如 $F/\mathbb Q_p$ 为非分歧二次扩张、$\pi=p$ 时，该标量是 $p^2$，不是 $p$。}

\DeclareTranslatorNote{13}{Kummer--$d\log$ 方块的符号}{按 2.2.9 对连接同态的约定，图中左方块实际反交换。为使它交换，可把左侧竖直箭头改为 $-d\log$；另一种表述是保留箭头而明确称其反交换。证明中令 $w=\log(\widetilde v^{p^m}h^{-1})$，已算得 $dw=-dh/h$，而 $\mu(\omega)$ 对 $dx=\omega$ 由余循环 $\sigma x-x$ 表示。因此证明最后一行的类应为 $-\mu_{\xi,n}(dh/h)$，不能省掉负号。此处的修正须与前后显式互反律中的符号一致。}

\DeclareTranslatorNote{14}{$\log\chi_G$ 的值域与整性}{一般应先把 $\log\chi_G$ 看作 $H^1(K,F)$ 中的元素；若需整系数形式，再乘一个足够可除的固定常数。分歧情形中，$\log:\mathcal O_F^\times\to\mathcal O_F$ 的整性断言不总成立。例如取 $F=K=\mathbb Q_3(\pi)$、$\pi^4=3$，并令 $v_\pi(\pi)=1$。则
\[
\log(1+\pi)=\sum_{j\ge1}(-1)^{j-1}\frac{\pi^j}{j}
\]
各项的赋值为 $j-4v_3(j)$，唯一最小值在 $j=3$ 处为 $-1$，故 $\log(1+\pi)\notin\mathcal O_F$。Lubin--Tate 特征可取到 $\mathcal O_F^\times$ 中这些值，因而需另行处理整性。这与先前译注指出的对数二次项正负号是两个不同的问题。}

\DeclareTranslatorNote{15}{de Rham 实现缺少滤过上标}{第二个同构应读为
\[
\Omega\otimes\mathbb Q_p\simeq D^0_{\mathrm{dR}}(\mathscr F),
\]
即取滤过的第 $0$ 部分。第一章 3.2.6 所定义的 $\Omega$ 是 $M^0_{\mathrm{dR}}$，第二章 2.1.3 中相应的 $\operatorname{coLie}$ 同构也以 $D^0_{\mathrm{dR}}$ 为源。在复乘情形，完整 de Rham 空间与该滤过部分的秩可能不同，不能省去上标。}

\DeclareTranslatorNote{16}{辅助映射 $a'$ 并非一般同构}{此处 $a'$ 上的 $\xrightarrow{\sim}$ 应改为普通箭头；经指定的系数特化后得到的 $a$ 才是所需同构。例如在有理数情形取 $r=1$、$L=\mathbb Q(i)$ 和奇二次特征，所写 $a'$ 对应实线性映射
\[
z\longmapsto(z\bmod\mathbb R,\bar z\bmod\mathbb R).
\]
定义域和值域的实维数均为 $2$，但其核为 $\mathbb R$、秩为 $1$。投影到所需的奇特征分量后，才得到一维同构。不能据原式把未经特化的 $a'$ 也当作同构。}

\DeclareTranslatorNote{17}{复乘局部计算应取 $\mathfrak p$ 分量}{这段局部计算应使用
\[
T_{\mathfrak p}E
=T_pE\otimes_{\mathcal O_K\otimes\mathbb Z_p}\mathcal O_{\mathfrak p},
\]
并把相应 Betti 格换成 $H_1(E(\mathbb C),\mathbb Z)\otimes_{\mathcal O_K}\mathcal O_{\mathfrak p}$。例如 $p$ 在复乘域中分裂时，$\mathcal O_{\mathfrak p}=\mathbb Z_p$，而完整的 $T_pE$ 的 $\mathbb Z_p$ 秩为 $2$，不可能以所写的一个元素作为 $\mathcal O_{\mathfrak p}$ 基；投影到 $\mathfrak p$ 分量后才有相应的秩一表述。}

\DeclareTranslatorNote{18}{复乘情形的系数代数}{这里应取
\[
A=K[\Gal(L/K)],
\]
而不是 $\mathbb Q[\Gal(L/K)]$。本段处于虚二次域情形，前文要求 $A$ 是 $K$ 的有限扩张的有限直积，并且使用 $K[\Gal(L/K)]$ 上的张量积。若 $\mathbb Q[G]$ 有所需的含幺 $K$ 代数结构，将其与增广商 $\mathbb Q[G]\to\mathbb Q$ 复合就会给出 $K\hookrightarrow\mathbb Q$，与 $K$ 为虚二次域矛盾。}

\DeclareTranslatorNote{19}{Siegel--Bombieri--Vaaler 定理的高度归一化}{定义 2.3(d) 沿用的是 Plücker 坐标向量的 sup 高度；定理 2.4 中不增加修正常数的版本，在无穷位须使用 Euclidean Plücker 范数。若保留当前 $h(A)$，可在定理右侧增加 $\tfrac12\log\binom NM$ 来吸收两种矩阵高度之差。例如 $k=\mathbb Q$、$A=(5,7,8)$ 时，当前 $h(A)=\log8$，但任何非零有理核向量化成本原整数向量后，最大坐标绝对值至少为 $3$：若各坐标绝对值至多为 $2$，由 $y\equiv z\pmod5$ 得 $y=z$、$x=-3y$，只能为零。因此任一核空间基的高度和至少为 $\log9>\log8$；$(3,-1,-1)$ 与 $(1,-3,2)$ 达到该下界。Euclidean Plücker 高度在此例中为 $\tfrac12\log138$，不会产生上述矛盾。}

\DeclareTranslatorNote{20}{Riemann--Roch 下界的一次项}{这里的 Riemann--Roch 公式还含有一次项：
\[
\chi(\mathcal O(dY))=\delta d^2-(g-1)(a_1+a_2)d+(g-1)^2.
\]
由 $h^2=0$ 只能直接推出 $h^0\ge\delta d^2-O(d)$；也可在 $d$ 充分大时使用 $h^0\ge(\delta/2)d^2$，不能省略低阶项而声称 $h^0\ge\delta d^2$。例如 $g=2,a_1=4,a_2=1,\delta=2$ 时，$Y=\Delta+3F_1$ 丰沛，对充分大且满足整除条件的 $d$，有 $h^0(\mathcal O(dY))=2d^2-5d+1<2d^2$。这一修正保留所需的二次增长，但须相应调整后续常数。}

\DeclareTranslatorNote{21}{单项式范围的严格端点}{两个范围应分别读为 $0\le b<N$ 和 $0\le b'<N$，上界不能取到 $N$。以一个因子说明：若取 $0\le a+b\le D$、$0\le b\le N$，当 $D$ 大时，单项式数为
\[
(N+1)(D+1)-N(N+1)/2,
\]
其一次项为 $(N+1)D$；目标空间 $H^0(C,\mathcal O(DNF))$ 的维数却为 $ND+1-g$。源空间维数已经大于目标，不能得到所声称的单射。改为严格不等式后，后面的线性方程构造也应使用相同的范围。}

\DeclareTranslatorNote{22}{范数比较所需的测度假设}{除书末附记已提到的 $X$ 紧致性外，还须对测度作非退化假设。为与给出的局部积分证明相配，可要求 $\nu$ 为光滑、处处正的体积测度。仅说“固定一个测度”不够：在紧致的 $X=\mathbb P^1$ 上取一点的 Dirac 测度 $\delta_P$，再取 $\mathcal O(1)$ 的一个在 $P$ 消失的非零截面 $s$，就有 $\|s\|_{L^2}=0$ 而 $\|s\|_{\sup}>0$，不可能有正常数 $c_{\mathscr L}$ 满足所写的不等式。}

\DeclareTranslatorNote{23}{Faltings 复形 (9.3) 末项的符号}{式 (9.3) 末项中的 $(s_i\operatorname{pr}_i-s_j\operatorname{pr}_j)^*L$ 应读为带加号的 $(s_i\operatorname{pr}_i+s_j\operatorname{pr}_j)^*L$。令
\[
D_\pm=\sum_{i<j}(s_i\operatorname{pr}_i\pm s_j\operatorname{pr}_j)^*L,
\qquad H=\sum_i s_i^2\operatorname{pr}_i^*L.
\]
则 $D_-+D_+=2(n-1)H$。源线丛为 $A=d(D_-+\delta H)$，中项为 $B=d(2n-2+\delta)H=A+dD_+$；乘子 $u_\alpha$ 属于 $\mathcal O(dD_+)$，故关系 $u_\alpha t_\beta-u_\beta t_\alpha$ 应落在 $\mathcal O(B+dD_+)$。再乘以 $dD_-$ 的截面才得到 (9.4)。这与 §11 中已另行记录的符号问题不是同一处。}

\DeclareTranslatorNote{24}{范数截面多次数中的总次数}{式 (13.10) 应读为
\[
e_i=d_i\prod_{j=1}^nN_j=d_iN,
\qquad N=\deg(\pi_1\times\cdots\times\pi_n).
\]
这里的原线丛是 $\pi^*M$，$M=\mathcal O(d_1,\ldots,d_n)$，由 $\Nm_\pi(\pi^*M)=M^N$，每个多次数都应乘总次数 $N$，不能漏掉 $N_i$。仅取 $n=1$，原式给出 $e_1=d_1$，而范数实际给出 $e_1=N_1d_1$。推论 18.3 证明中再次写出的 $e_i=d_iN/N_i$ 也应同步按 $d_iN$ 理解。}

\DeclareTranslatorNote{25}{命题 14.1 中 $\Gamma_1$ 的 $\epsilon$ 符号}{这里应为
\[
\Gamma_1=\Gamma\!\left(\mathscr W,d(2n-2-\epsilon)\sum_i s_i^2\operatorname{pr}_i^*L\right)^a.
\]
所用复形来自把 $\delta=-\epsilon$ 代入 (9.4)，所以 $2n-2+\delta$ 应成为 $2n-2-\epsilon$。保留加号会使构造属于另一条线丛，无法给出命题所需的 $L_{-\epsilon,s}$ 的截面。}

\DeclareTranslatorNote{26}{非阿基米德 Hensel 引理的完备性}{此处应增加 $K$ 关于给定非阿基米德绝对值完备的假设，才可从 Newton 序列推出它在 $K$ 中收敛到根。例如 $K=\mathbb Q$ 配 $7$ 进绝对值，取 $f(T)=T^2-2$、$\alpha_0=3$，则 $|f(3)|_7=1/7<1=|f'(3)|_7^2$，满足所写条件，但根 $\sqrt2$ 不在 $\mathbb Q$ 中。后文实际使用的完备域 $\mathbb C_v$ 满足补充后的假设。}

\DeclareTranslatorNote{27}{复数 Newton 引理应假设全纯}{此处应要求各 $f_i$ 全纯。若只要求实意义下的 $C^2$，不能仅由 $m\times m$ 复 Jacobian 和 $\partial^2/\partial z_i\partial z_j$ 的估计推出结论。例如 $m=1$、$f(z)=z+\bar z+i/4$、$\alpha_0=0$、$B=1$ 时，$f_z=1$、$f_{zz}=0$，满足所写估计，但 $\operatorname{Im}f(z)=1/4$，根不存在。若要保留一般 $C^2$ 映射，应改用完整的 $2m\times2m$ 实 Jacobian 及相应 Hessian，并重新表述估计。}

\DeclareTranslatorNote{28}{乘积定理高度界所需的整性}{定理 18.1 及推论 18.3 的 (18.2)、(18.4)，只用阿基米德范数控制高度时，须补上数域及整性条件：截面应来自固定整数模型，且所用模型、格及度量须在估计中固定。另一种修正是把右侧改为计入有限位分母贡献的适当射影高度。否则把任意截面 $\gamma$ 换成 $\gamma/M$，零集、指标和次数不变，阿基米德对数范数却可趋于负无穷。例如在定理 18.1 中取 $n=1$、$\Pi=\mathbb P^1_{\mathbb Q}$、$e_1=1$、$\gamma'=X_0/M$、$x=[0:1]$，并令 $0<\epsilon_3\le1$；所需真子簇只能是固定点 $\{x\}$，高度固定，但 (18.2) 右侧随 $M\to\infty$ 趋于负无穷。因此不能对任意有理截面直接使用现有的高度界。}

\DeclareTranslatorNote{29}{范数估计 (18.5) 的总次数幂}{若有限映射总次数为 $N$ 且 $\gamma'=\Nm(\gamma)$，右侧第一个因子应为
\[
\left(\prod_{v\mid\infty}\|\gamma\|_{\sup,v}\right)^N.
\]
因此 (18.5) 应在该因子后保留原有的指数误差项。缩放可见这一幂次不可省略：把 $\gamma$ 换成 $m\gamma$，范数截面变为 $m^N\gamma'$；原左侧按 $m^N$ 增长，原右侧却只按 $m$ 增长，$N>1$ 时矛盾。后文取对数并使用 (18.5) 的估计也须同步计入这个 $N$；由于相关次数有界，吸收常数时仍须明确这一来源。}


\begin{document}
\frontmatter
\begin{titlepage}
\thispagestyle{empty}
\begin{tikzpicture}[remember picture,overlay]
\fill[Ink] (current page.north west) rectangle ([xshift=13mm]current page.south west);
\draw[Accent,line width=1.1pt] ([xshift=26mm,yshift=-35mm]current page.north west)--([xshift=-22mm,yshift=-35mm]current page.north east);
\node[anchor=south east,rotate=90,text=white,font=\sffamily\small] at ([xshift=9mm,yshift=22mm]current page.south west) {LECTURE NOTES IN MATHEMATICS \quad 1553};
\draw[Ink!12,line width=0.5pt] ([xshift=-60mm,yshift=56mm]current page.south east) circle (23mm);
\draw[Ink!12,line width=0.5pt] ([xshift=-60mm,yshift=56mm]current page.south east) ellipse (36mm and 15mm);
\draw[Ink!12,line width=0.5pt] ([xshift=-60mm,yshift=56mm]current page.south east) ellipse (15mm and 36mm);
\end{tikzpicture}
\vspace*{28mm}
{\sffamily\small\color{Accent}C.I.M.E.\quad 特伦托，1991}\par
\vspace{16mm}
{\fontsize{35}{46}\selectfont\bfseries\color{Ink}算术代数几何}\par
\vspace{8mm}
{\Large\color{Ink}Arithmetic Algebraic Geometry}\par
\vspace{10mm}
{\large 1991 年 C.I.M.E. 第二期暑期学校讲义}\par
\vspace{16mm}
{\large Jean-Louis Colliot-Thélène\\[5pt]Kazuya Kato\\[5pt]Paul Vojta}\par
\vspace{9mm}
{\color{Muted}Edoardo Ballico\quad 编}\par
\vfill
{\color{Muted}中文译本·译者补注版\quad ·\quad XeLaTeX 排印}\par
\vspace{5mm}
{\small\color{Muted}原版：Springer-Verlag，1993\quad｜\quad 数学讲义丛书第 1553 卷}\par
\vspace{15mm}
\end{titlepage}

\chapter*{原书出版信息}
\addcontentsline{toc}{chapter}{原书出版信息}
\noindent\textbf{数学讲义丛书（Lecture Notes in Mathematics）1553}

丛书编委：A. Dold（海德堡）、B. Eckmann（苏黎世）、F. Takens（格罗宁根）。

佛罗伦萨 C.I.M.E. 基金会分系列；顾问：Roberto Conti。

本书收录 1991 年 6 月 24 日至 7 月 2 日在意大利特伦托举行的国际数学暑期中心（Centro Internazionale Matematico Estivo，C.I.M.E.）第二期讲习班讲义。

\subsection*{作者与编者}
\begin{description}
\item[Jean-Louis Colliot-Thélène]巴黎南大学数学系，425 号楼，F-91005 Orsay Cedex，法国。
\item[Kazuya Kato（加藤和也）]东京工业大学数学系，Oh-Okayama，Meguro-ku，Tokyo，日本。
\item[Paul Vojta]加利福尼亚大学伯克利分校数学系，Berkeley，CA 94720，美国。
\item[Edoardo Ballico（编）]特伦托大学数学系，38050 Povo，Trento，意大利。
\end{description}

\noindent 1991 年数学主题分类：14C15，14C25，14C35，19E15，14G40，11J99，11M99，11G99，11R23，14M20。

\noindent ISBN 3-540-57110-8\quad Springer-Verlag，柏林、海德堡、纽约。\\
ISBN 0-387-57110-8\quad Springer-Verlag，纽约、柏林、海德堡。

\subsection*{原书版权声明译文}
本书受版权保护。全部或部分内容的各项权利均予保留，特别包括翻译、重印、插图再利用、朗诵、广播、缩微胶片或其他形式的复制，以及在数据库中储存等权利。只有在 1965 年 9 月 9 日《德国著作权法》现行条文允许的范围内，才可复制本书或其中部分内容；使用时始终须取得 Springer-Verlag 的许可。违反规定者将依德国著作权法受到追究。

\noindent © Springer-Verlag Berlin Heidelberg 1993。德国印刷。\\
印刷与装订：Druckhaus Beltz，Hemsbach/Bergstr.\\
46/3140-543210；使用无酸纸印刷。

原文件另附 2008 年版勘误页，其 ISBN 为 978-3-540-57110-0，版权标识为 © Springer-Verlag Berlin Heidelberg 2008。该勘误及补充文献完整译于书末。

\chapter*{译本使用说明}
\addcontentsline{toc}{chapter}{译本使用说明}
本译本按原书三篇讲义的顺序编排，保留定理、命题、引理、公式及文献的原有编号。页边“原书”标记对应原书印刷页码；译本页码独立编排。文中叙述及“最近”“尚未发表”等时间表述以原书写作年代为准。

技术名词采用“中文（原文术语）”的形式：第一篇括注法语，后两篇括注英语。单击术语可跳至文末中英法对照表；表中“正文位置”列的页码可返回相应出现位置。同页有多次出现时，重复页码按出现先后排列，分别指向各处。译者补充的第三语言对应词仅供对照。

原文件包含两页无可提取文字层的后补勘误，译本将其单列于书末，并在受实质性修正影响的正文处加注。原书参考文献的作者名、题名和出版信息保留原文，便于检索；参与者名单中的人名与通信地址也予以保留。

本次补注版依据另行收到的勘误说明，在相应正文处加入编号 1--29 的“译者补注”，并附可往返跳转的索引。补注保留该说明提出的修正式、必要假设和理由；定理 6.4 的问题明确记为证明范围的缺口。原书表述及作者正式勘误保留，供读者对照。

\chapter*{序言}
\addcontentsline{toc}{chapter}{序言}
1991 年 6 月 24 日至 7 月 2 日，C.I.M.E. 举办的“算术代数几何”讲习班在意大利特伦托的 Villa Madruzzo 举行。

讲习班安排了以下系列讲座：
\begin{itemize}
\item Jean-Louis Colliot-Thélène：挠代数闭链与代数 $K$ 理论；
\item Kazuya Kato：经由 $B_{\dR}$ 研究 Hasse--Weil $L$ 函数的岩泽理论方法；
\item Christophe Soulé：Arakelov 几何；
\item Paul Vojta：算术几何在丢番图逼近中的应用。
\end{itemize}
此外，参加者还作了以下专题报告：
\begin{itemize}
\item Dan Abramovich：阿贝尔簇及 Jacobian 簇的子簇；
\item Luca Barbieri Viale：通过上同调理论研究双有理不变量；
\item Torsten Ekedahl：a）中国剩余定理的一个无限版本；b）具有一般分歧类型的扩张的密度；
\item Frans Oort：a）阿贝尔簇的复乘提升；b）Coleman 关于复数域上具有复乘的 Jacobian 簇的一个猜想；c）Newton 多边形与阿贝尔簇；
\item Angelo Vistoli：双变相交理论与 Alexander 对偶；
\item Christoph Wirsching：$G(2,4)$ 上的 Quillen 度量。
\end{itemize}
本卷收入其中三个系列讲座的扩充稿。Soulé 讲座的内容可参见 Ch. Soulé、D. Abramovich、J. F. Burnol 与 J. K. Kramer 最近出版的专著 \emph{Arakelov Geometry}，Cambridge Studies in Advanced Mathematics 33，Cambridge University Press，1992。

四位主讲人都煞费苦心，从有志钻研而尚非本领域专家的研究生所能理解的层次讲起，逐步引向研究前沿，乃至前沿之外。在本次 C.I.M.E. 讲习班之前、期间和之后，他们还以各种方式为听众提供了进一步的帮助。仅举一例：J.-L. Colliot-Thélène 早在 1990 年便编出一份长篇参考书目，既含入门资料，也含高深文献，并附有详细的内容说明；这份书目曾广泛流传。收入本书的三篇讲义，其作者在撰写时也付出了同样的心力。

在会议筹备和本书编印过程中，许多人在数学和具体事务两方面都给予了宝贵帮助，包括所有讲授者和参加者。我谨向他们致以特别的谢意。

\begin{flushright}Edoardo Ballico\end{flushright}

\begingroup\termanchorsfalse\tableofcontents\endgroup
\mainmatter
\input{generated/01-colliot-a}
\input{generated/02-colliot-b}
\input{generated/03-kato-a}
\input{generated/04-kato-b}
\input{generated/05-kato-c}
\input{generated/06-kato-d}
\input{generated/07-vojta-a}
\input{generated/08-vojta-b}
\backmatter
\chapter*{第一篇勘误与评论}
\addcontentsline{toc}{chapter}{第一篇勘误与评论}
\phantomsection\label{errata-ct}
\noindent\textcolor{Muted}{以下译自原文件附载的 2008 年勘误两页，对应第一篇原书第 1--49 页。所引页码及行号均指原书；负行号从页底向上数。}

\begin{description}
\item[第 14 页，页底]
Chad Schoen（\emph{J. Algebraic Geom.}\ \textbf{11}（2002），41--100）已经证明，这个映射一般并非单射。
\hfill\hyperref[orig:19]{返回正文}

\item[第 18 页，倒数第 5 行]
应为“réduction”（约化）。\hfill\hyperref[orig:23]{返回正文}

\item[第 24 页，第 19 行]
应为“D'après”（根据）。\hfill\hyperref[orig:29]{返回正文}

\item[第 32 页，第 2 行]
应为“apparue”（出现）。\hfill\hyperref[orig:37]{返回正文}

\item[第 38 页，第 24 行]
应引用 [CR3]。\hfill\hyperref[orig:43]{返回正文}

\item[第 40 页，定理 8.4 及备注 8.4.1]
定理 8.4 中关于 Albanese 簇具有好约化的假设不能删去；R. Parimala 与 V. Suresh 的一个例子说明了这一点，见 \emph{Inventiones math.}\ \textbf{122}（1995），83--117。
\hfill\hyperref[orig:45]{返回正文}

\item[第 41 页，第 6 行]
应读作
\[
 H^2\!\left(K,H^1_{\et}\!\left(\overline X,\mathbb Q/\mathbb Z(2)\right)^0\right).
\]
\hfill\hyperref[orig:46]{返回正文}

\item[第 42 页，倒数第 5 行至倒数第 1 行]
在定理 8.6 中补充假设：$X/R$ 的相对维数为 $2$。
\hfill\hyperref[orig:47]{返回正文}

\item[第 45 页，第 10 行]
删去一个“le”。\hfill\hyperref[orig:50]{返回正文}

\item[第 45 页，倒数第 2 行]
补充：利用从 $L$ 到 $k$ 的迹，可以看出映射
\[
 H^1(X,\mathcal K_2)\otimes\mathbb Q
 \longrightarrow
 \bigoplus_{p\in\Spec A^{(1)}}\Pic(\mathbb X_p)\otimes\mathbb Q
\]
是满射。\hfill\hyperref[orig:50]{返回正文}

\item[第 46 页，第 2 行]
将
\[
 CH_2(\mathbb X_B)_{\tors}\longrightarrow CH^2(X_L)_{\tors}
\]
替换为
\[
 CH_2(\mathbb X)_{\tors}\longrightarrow CH^2(X)_{\tors}.
\]
勘误页的左端仍印为下标 $2$；正文依据余维二的上下文采用上标 $2$，并另加译注。\hfill\hyperref[orig:51]{返回正文}

\item[第 46 页，倒数第 14 行]
应为“exacte”（正合）。\hfill\hyperref[orig:51]{返回正文}

\item[第 46 页，倒数第 12 行]
将 $H^i(\mathbb X,\mathbb Z_l(j))$ 换成该群模去其挠子群所得的商群。
\hfill\hyperref[orig:51]{返回正文}

\item[第 46 页，倒数第 6 行]
删去“les”。\hfill\hyperref[orig:51]{返回正文}

\item[第 48 页，参考文献 [Sb2]]
Per Salberger 的文章已经发表：\emph{Torsion cycles of codimension 2 and $l$-adic realizations of motivic cohomology}，载于 \emph{Séminaire de théorie des nombres, Paris 1991--1992}（Sinnou David 编），Progress in Mathematics \textbf{116}，Birkhäuser（1993），第 247--277 页。
\hfill\hyperref[orig:53]{返回正文}
\end{description}

\section*{补充参考文献}
\begin{enumerate}
\item J.-L. Colliot-Thélène，\emph{L'arithmétique des zéro-cycles}（1993 年 9 月波尔多算术研讨会报告），\emph{Journal de théorie des nombres de Bordeaux}\ \textbf{7}（1995），51--73。题名译为“零闭链的算术”。
\item J.-L. Colliot-Thélène，\emph{Zéro-cycles sur les surfaces sur un corps $p$-adique : quelques résultats obtenus au moyen de la K-théorie algébrique}，为 Morelia（2003 年 6--7 月）和 Sestri Levante（2004 年 6--7 月）会议准备的讲义，可在原勘误所列网页 \url{http://www.math.u-psud.fr/~colliot/liste-cours-exposes.html} 获取。题名译为“$p$ 进域上曲面的零闭链：借助代数 $K$ 理论得到的若干结果”。
\item Masanori Asakura 与 Shuji Saito，\emph{Surfaces over a $p$-adic field with infinite torsion in the Chow group of $0$-cycles}，\emph{Algebra and Number Theory}\ \textbf{1}（2007），163--181。题名译为“零闭链 Chow 群中含有无限挠子群的 $p$ 进域上曲面”。
\item Shuji Saito 与 Kanetomo Sato（附 Uwe Jannsen 撰写的附录），\emph{A finiteness theorem on zero-cycles over $p$-adic fields}，待刊于 \emph{Annals of Math.}。题名译为“$p$ 进域上零闭链的一个有限性定理”。此处“待刊”沿用勘误的记载。
\item Shuji Saito 与 Kanetomo Sato，\emph{Torsion cycle class maps in codimension two of arithmetic schemes}，预印本 arXiv:math/0612081。题名译为“算术概形余维二的挠闭链类映射”。
\end{enumerate}

\chapter*{译读附记}
\addcontentsline{toc}{chapter}{译读附记}
以下说明补充记录第三篇中沿用的原书记号及表述。页码均指原书，右侧链接返回相应正文。第一、二篇的实质性校订及疑点已在相关页的译注中说明；第一篇的正式勘误另见前一附录。

\begin{description}
\item[第 180 页，延拓 $F$ 的叙述]
上文取 $F=K_C/(2g-2)$，这里延拓时直接写 $\omega_{X/B}$；译文沿用这一记号，未另添归一化因子。
\hfill\hyperref[orig:185]{返回正文}
\item[第 181 页，引理 8.1]
陈述中未显式假设 $X$ 紧致，证明中则使用了紧致性。本次补充了测度非退化条件的说明，见\hyperlink{translator-note-22}{译者补注22}。\hfill\hyperref[orig:186]{返回正文}
\item[第 182 页，引理 8.2]
第二项的分母为 $2d_1$，紧接的文字则给出 $(2g-1)/r$；两处均沿用原书。\hfill\hyperref[orig:187]{返回正文}
\item[第 183 页，正合列附近]
$a,b$ 既用于整除关系 $ba^2=d$，也用于直和重数；译文沿用原书记号。\hfill\hyperref[orig:188]{返回正文}
\item[第 187 页，推论 11.4]
原书定义 $\mathscr P_{ij}=(\operatorname{pr}_i+\operatorname{pr}_j)^*L-\operatorname{pr}_i^*L-\operatorname{pr}_j^*L$，后面对 $(s_i\operatorname{pr}_i-s_j\operatorname{pr}_j)^*L$ 的展开末项仍印为 $+s_is_j\mathscr P_{ij}$；译文保留该加号。\hfill\hyperref[orig:192]{返回正文}
\item[第 189 页，引理 12.1(a)]
文字称常数 $c_1,c_2$，公式右端却使用 $c_3^{\delta}$；译文保留原书的对应字母。\hfill\hyperref[orig:194]{返回正文}
\item[第 193 页，式 (13.11)]
指数中 $c\sum d_i h(X_i)+c''\sum d_i$ 的加号沿用原书。\hfill\hyperref[orig:198]{返回正文}
\item[第 196 页，Hensel 引理的证明]
原书 Taylor 展开的一个式子写为 $f(\alpha')=f(\alpha)-J_f(\alpha)(\alpha'-\alpha)+\cdots$，一次项前的减号予以保留。\hfill\hyperref[orig:201]{返回正文}
\item[第 197 页，引理 15.8]
原文称 $J_f$ 为 Jacobian determinant，但定义和后文使用的是前述矩阵；译文写作“如引理 15.3 中的 Jacobian”。\hfill\hyperref[orig:202]{返回正文}
\item[第 198 页，归纳步骤 (i)]
求和下限原印为 $j=1$，后文 (15.11) 的对应求和则从 $0$ 开始；两处沿用原书。\hfill\hyperref[orig:203]{返回正文}
\item[第 200 页，参考文献引用]
正文写 [So 1]，文末相应文献项标为 [So]；译文保留原有引用键。\hfill\hyperref[orig:205]{返回正文}
\item[第 201 页，引理 16.2 的证明]
原书一处列 $f_1,\ldots,f_{m_i}$，前文则为 $f_1,\ldots,f_{M-m_i}$；译文保留两处原有范围。\hfill\hyperref[orig:206]{返回正文}
\item[第 204 页，第 17 节]
原书所写 $\deg\bigl((s_i^2\operatorname{pr}_i-s_j^2\operatorname{pr}_j)^*L\bigr)|_E$ 的态射系数含平方，紧接右端则含 $s_i^2,s_j^2$；译文保留这一原式。\hfill\hyperref[orig:209]{返回正文}
\end{description}

\chapter*{本次译者补注索引}
\addcontentsline{toc}{chapter}{本次译者补注索引}
本次补注依据 2026 年 11 月 5 日收到的勘误说明编写，编号 1--29 与该说明一一对应。各注列出说明所建议的修正式、适用条件及理由；第 2 项单列证明范围的缺口。作者正式勘误、原有译注与本次补注分别保留。

本版正文保留原书表述，相关位置以“译者补注”说明应如何理解或修正；同一问题再次出现时，另附跳转提示。这里使用章节和公式号定位，并列出本版实际页码。单击条目标题可跳至补注，单击注内的补注编号可返回本索引。

\begin{enumerate}[label=\arabic*.,leftmargin=2.5em,itemsep=7pt]
\item \hypertarget{translator-index-1}{}\hyperlink{translator-note-1}{\csuse{translator-note-title-1}}。
第一篇，命题1.1后的Picard下降论证；第\,\pageref*{translator-note:1}\,页。
\item \hypertarget{translator-index-2}{}\hyperlink{translator-note-2}{\csuse{translator-note-title-2}}。
第一篇，定理6.4及其证明；第\,\pageref*{translator-note:2}\,页。
\item \hypertarget{translator-index-3}{}\hyperlink{translator-note-3}{\csuse{translator-note-title-3}}。
第一篇，备注8.4.2；第\,\pageref*{translator-note:3}\,页。
\item \hypertarget{translator-index-4}{}\hyperlink{translator-note-4}{\csuse{translator-note-title-4}}。
第一篇，定理9.1证明的第二个交换图及说明；第\,\pageref*{translator-note:4}\,页。
\item \hypertarget{translator-index-5}{}\hyperlink{translator-note-5}{\csuse{translator-note-title-5}}。
第二篇，I §2.2.3(3)；第\,\pageref*{translator-note:5}\,页。
\item \hypertarget{translator-index-6}{}\hyperlink{translator-note-6}{\csuse{translator-note-title-6}}。
第二篇，I §2.2.4；第\,\pageref*{translator-note:6}\,页。
\item \hypertarget{translator-index-7}{}\hyperlink{translator-note-7}{\csuse{translator-note-title-7}}。
第二篇，I 猜想2.2.7(ii)及§2.2.6(4)；第\,\pageref*{translator-note:7}\,页。
\item \hypertarget{translator-index-8}{}\hyperlink{translator-note-8}{\csuse{translator-note-title-8}}。
第二篇，I §2.2.9；第\,\pageref*{translator-note:8}\,页。
\item \hypertarget{translator-index-9}{}\hyperlink{translator-note-9}{\csuse{translator-note-title-9}}。
第二篇，I §3.3.5；第\,\pageref*{translator-note:9}\,页。
\item \hypertarget{translator-index-10}{}\hyperlink{translator-note-10}{\csuse{translator-note-title-10}}。
第二篇，I 引理3.5.5；第\,\pageref*{translator-note:10}\,页。
\item \hypertarget{translator-index-11}{}\hyperlink{translator-note-11}{\csuse{translator-note-title-11}}。
第二篇，II §2.1.4；第\,\pageref*{translator-note:11}\,页。
\item \hypertarget{translator-index-12}{}\hyperlink{translator-note-12}{\csuse{translator-note-title-12}}。
第二篇，II §2.2.5；第\,\pageref*{translator-note:12}\,页。
\item \hypertarget{translator-index-13}{}\hyperlink{translator-note-13}{\csuse{translator-note-title-13}}。
第二篇，II 引理2.2.8及§2.2.10证明末行；第\,\pageref*{translator-note:13}\,页。
\item \hypertarget{translator-index-14}{}\hyperlink{translator-note-14}{\csuse{translator-note-title-14}}。
第二篇，II §2.3.1；第\,\pageref*{translator-note:14}\,页。
\item \hypertarget{translator-index-15}{}\hyperlink{translator-note-15}{\csuse{translator-note-title-15}}。
第二篇，III §1.2.5的de Rham同构；第\,\pageref*{translator-note:15}\,页。
\item \hypertarget{translator-index-16}{}\hyperlink{translator-note-16}{\csuse{translator-note-title-16}}。
第二篇，III §1.2.5的辅助映射$a'$；第\,\pageref*{translator-note:16}\,页。
\item \hypertarget{translator-index-17}{}\hyperlink{translator-note-17}{\csuse{translator-note-title-17}}。
第二篇，III 定理1.2.6证明中的局部计算；第\,\pageref*{translator-note:17}\,页。
\item \hypertarget{translator-index-18}{}\hyperlink{translator-note-18}{\csuse{translator-note-title-18}}。
第二篇，III 定理1.2.6证明中的系数代数；第\,\pageref*{translator-note:18}\,页。
\item \hypertarget{translator-index-19}{}\hyperlink{translator-note-19}{\csuse{translator-note-title-19}}。
第三篇，定义2.3(d)及定理2.4；第\,\pageref*{translator-note:19}\,页。
\item \hypertarget{translator-index-20}{}\hyperlink{translator-note-20}{\csuse{translator-note-title-20}}。
第三篇，§7第一阶段的Riemann--Roch下界；第\,\pageref*{translator-note:20}\,页。
\item \hypertarget{translator-index-21}{}\hyperlink{translator-note-21}{\csuse{translator-note-title-21}}。
第三篇，§7第二阶段的单项式范围；第\,\pageref*{translator-note:21}\,页。
\item \hypertarget{translator-index-22}{}\hyperlink{translator-note-22}{\csuse{translator-note-title-22}}。
第三篇，引理8.1；第\,\pageref*{translator-note:22}\,页。
\item \hypertarget{translator-index-23}{}\hyperlink{translator-note-23}{\csuse{translator-note-title-23}}。
第三篇，Faltings复形(9.3)；第\,\pageref*{translator-note:23}\,页。
\item \hypertarget{translator-index-24}{}\hyperlink{translator-note-24}{\csuse{translator-note-title-24}}。
第三篇，式(13.10)及推论18.3证明；第\,\pageref*{translator-note:24}\,页。
\item \hypertarget{translator-index-25}{}\hyperlink{translator-note-25}{\csuse{translator-note-title-25}}。
第三篇，命题14.1证明中的$\Gamma_1$；第\,\pageref*{translator-note:25}\,页。
\item \hypertarget{translator-index-26}{}\hyperlink{translator-note-26}{\csuse{translator-note-title-26}}。
第三篇，引理15.3及其前置设定；第\,\pageref*{translator-note:26}\,页。
\item \hypertarget{translator-index-27}{}\hyperlink{translator-note-27}{\csuse{translator-note-title-27}}。
第三篇，引理15.8；第\,\pageref*{translator-note:27}\,页。
\item \hypertarget{translator-index-28}{}\hyperlink{translator-note-28}{\csuse{translator-note-title-28}}。
第三篇，定理18.1、推论18.3及式(18.2)、(18.4)；第\,\pageref*{translator-note:28}\,页。
\item \hypertarget{translator-index-29}{}\hyperlink{translator-note-29}{\csuse{translator-note-title-29}}。
第三篇，范数估计(18.5)及后续对数估计；第\,\pageref*{translator-note:29}\,页。
\end{enumerate}

\chapter*{讲习班参加者名单}
\addcontentsline{toc}{chapter}{讲习班参加者名单}
\noindent C.I.M.E.“算术代数几何”讲习班。以下保留原书姓名及历史通信地址；机构名称作中文翻译。

\origpage{216}{209}
\begin{description}[font=\normalfont\bfseries,style=nextline,leftmargin=0pt,labelwidth=0pt,labelsep=0pt,itemindent=0pt,itemsep=7pt]
\item[D. Abramovich]哈佛大学数学系，One Oxford Str., Cambridge, MA 02138。
\item[M. Andreatta]Via Anzoletti 14, 38100 Trento。
\item[L. Barbieri Viale]数学系，Via L.B. Alberti 4, 16132 Genova。
\item[A.-S. Basarab]罗马尼亚科学院数学研究所，Str. Academiei 14, 70700 Bucharest。
\item[M. Berni]Via del Cuore 6, 56100 Pisa。
\item[G. Canuto]数学系，Strada Nuova 65, 27100 Pavia。
\item[D. Chirica]罗马尼亚科学院数学研究所，Str. Academiei 14, 70700 Bucharest。
\item[J. Coanda]特伦托大学数学系，38050 Povo, Trento。
\item[C.I. Cobeli]罗马尼亚科学院数学研究所，Str. Academiei 14, 70700 Bucharest。
\item[C.-S. Dalawat]巴黎南大学数学系，F-91405 Orsay Cedex。
\item[M.A. de Cataldo]圣母大学数学系，Notre Dame, USA。
\item[T. Ekedahl]斯德哥尔摩大学数学系，S-11385 Stockholm。
\item[G. Elencwajg]尼斯大学数学实验室，Parc Valrose, F-06034 Nice Cedex。
\item[M. Flexor]巴黎第九大学，Orsay 91, Paris。
\item[N. Gavioli]特伦托大学数学系，38050 Povo, Trento。
\item[W. Gubler]苏黎世联邦理工学院数学系，8092 Zürich。
\item[K. Ha Huy]国际理论物理中心（ICTP）数学部，Box 586, 34100 Trieste。
\item[L. Lafforgue]10 rue Louis Gaudry, 92 160 Antony, France。
\item[F. Lecomte]Louis Pasteur 大学数学系，7 rue René Descartes, 67084 Strasbourg Cedex。
\item[V. Monti]特伦托大学数学系，38050 Povo, Trento。
\item[M. Nakamaye]耶鲁大学数学系，Box 2155, Yale Station, New Haven, CT 06520。
\item[L. Narváez-Macarro]数学学院，Tarfia s/n, 41020 Sevilla。
\item[F. Oort]数学研究所，Budapestlaan 6, Utrecht, NL。
\item[C. Pedrini]数学系，Via L.B. Alberti 4, 16132 Genova。
\item[D. Portelli]数学科学系，Piazzale Europa 1, 34127 Trieste。
\item[A. Previtali]Via Sant'Eurosia 15, 22064 Casatenovo, Como。
\item[Q.-V. Pham]Fehlinghöhe 21, 2000 Hamburg 60。
\end{description}
\origpage{217}{210}
\begin{description}[font=\normalfont\bfseries,style=nextline,leftmargin=0pt,labelwidth=0pt,labelsep=0pt,itemindent=0pt,itemsep=7pt]
\item[B. Russo]Via V. Veneto 821E, Bolzano。
\item[R. Salvati Manni]“La Sapienza”大学数学系，P.le A. Moro 2, 00185 Roma。
\item[E. Schiavi]特伦托大学数学系，38050 Povo, Trento。
\item[R. Schoof]特伦托大学数学系，38050 Povo, Trento。
\item[C. Soulé]高等科学研究所（IHES），35 Route de Chartres, F-91440 Bures-sur-Yvette。
\item[A. Thorup]数学研究所，Universitetsparken 5, DK-2100 København Ø。
\item[A. Vistoli]Via Irma Bandiera 121, 40024 Crevalcore, Bologna。
\item[C. Wirsching]慕尼黑大学数学研究所，Theresienstr. 39, D-8000 München 2。
\end{description}

\input{generated/glossary}
\end{document}
